\hypertarget{chapter-3-document-architecture-not-document-syntax}{%
\chapter{Document Architecture, Not Document Syntax}\label{chapter-3-document-architecture-not-document-syntax}}

Most LaTeX instruction stops at the level of syntax: how to write a section heading, how to insert a figure, how to format a table. That level of instruction is necessary but not sufficient for anyone working on a document longer than a handful of pages, because syntax says nothing about how to organize a project once it grows past the point where a single file is still convenient to write, review, and rebuild. This chapter is about the decisions that come before any of that syntax is written: how a book-length project is laid out on disk, how its parts relate to one another, and how that structure is kept consistent as the project grows over months or years rather than days.

\hypertarget{designing-a-project-before-writing-a-sentence}{%
\section{Designing a Project Before Writing a Sentence}\label{designing-a-project-before-writing-a-sentence}}

A short paper can be a single \texttt{.tex} file without cost. A book cannot, and the reasons are practical rather than aesthetic. A single-file monograph becomes slow to compile, since every edit forces the entire document to be reprocessed; it becomes difficult to review, since a diff against the previous version mixes changes from unrelated chapters together; and it becomes difficult to reason about, since finding the source of a given passage means searching rather than navigating. Treating chapters, appendices, figures, and the bibliography as separate concerns from the outset avoids all three problems, and doing so before the first chapter is drafted is considerably easier than retrofitting structure onto a project that has already grown large as a single file.

The separation that matters most is between content and assembly. Each chapter should live in its own file, containing only that chapter's prose, definitions, and figures; a single top-level file should contain the document class, the preamble, and a sequence of \texttt{\textbackslash{}input} or \texttt{\textbackslash{}include} commands that assembles the chapters into the finished book. This means a chapter file is never compiled in isolation without the correct preamble, but it also means a chapter can be edited, reviewed, and reasoned about as a self-contained unit, and that the top-level file serves as a table of contents for the source itself, not only for the rendered output.

\hypertarget{input-versus-include}{%
\section{\texorpdfstring{\texttt{\textbackslash{}input} versus \texttt{\textbackslash{}include}}{\textbackslash input versus \textbackslash include}}\label{input-versus-include}}

LaTeX offers two commands for pulling one file into another, and the difference between them matters for a project of any real size. \texttt{\textbackslash{}input} performs a direct textual inclusion: the named file's contents are read in place, with no side effects beyond whatever the file itself does. \texttt{\textbackslash{}include} does more: it forces a page break before and after the included file, and, critically, it interacts with \texttt{\textbackslash{}includeonly} to allow partial compilation, where only the listed files are actually typeset while the others are skipped, with their page numbers and cross-references still available from the existing auxiliary files.

For a book-length project under active revision, this distinction is not a minor convenience. Recompiling a five-hundred-page manuscript to check a single paragraph's line breaks is slow enough to disrupt a working rhythm; recompiling only the chapter currently being edited, while still resolving references into the rest of the book correctly, is not. The tradeoff is that \texttt{\textbackslash{}include} forces page breaks between included files, which is appropriate at chapter boundaries but wrong for finer-grained splitting, such as separating a chapter's figures or a long proof into their own file purely for editing convenience. A common and effective convention is to use \texttt{\textbackslash{}include} at the chapter level and \texttt{\textbackslash{}input} for everything finer, so that partial compilation aligns with chapter boundaries while sub-chapter organization stays flexible.

\hypertarget{directory-conventions-for-a-large-project}{%
\section{Directory Conventions for a Large Project}\label{directory-conventions-for-a-large-project}}

Once a project spans dozens of chapter files, generated figures, a bibliography, and assorted build artifacts, an unstructured flat directory becomes as much of an obstacle as an unstructured single file was. A directory layout that separates these categories from the start pays for itself repeatedly: a \texttt{chapters/} directory holding one file per chapter, numbered or named so their order is legible without opening them; a \texttt{figures/} directory holding source images and, separately, a location for figures generated by an external script or program, so that regenerated output does not get confused with hand-authored artwork; a \texttt{bib/} directory holding the bibliography database or the \texttt{thebibliography} source, kept as its own file even when the project is not using BibTeX, so that citation data is not duplicated across chapters; and a build directory, excluded from version control, holding the \texttt{.aux}, \texttt{.log}, \texttt{.toc}, and other files a compiler regenerates on every run.

This layout does more than keep the project tidy. It makes the boundary between what is authored and what is generated explicit, which matters once a project starts pulling figures from simulation output or data pipelines, a workflow covered later in this book. A generated figure that lives in a directory clearly marked as generated can be deleted and rebuilt without anxiety; a generated figure indistinguishable from a hand-drawn one cannot, and the distinction is easy to lose once a project has been worked on intermittently for a year or more.

\hypertarget{consistency-across-files}{%
\section{Consistency Across Files}\label{consistency-across-files}}

Splitting a document into many files introduces a problem a single file never has: keeping definitions, notation, and terminology consistent across pieces that are edited separately, sometimes months apart. A macro defined only in the preamble of the top-level file is available everywhere, which is the right home for anything used across more than one chapter; a macro defined inside a chapter file and needed elsewhere is a sign that it belongs in a shared preamble file instead, included once and referenced by every chapter that needs it.

Cross-references are the sharpest version of this problem. A \texttt{\textbackslash{}ref} or \texttt{\textbackslash{}cite} in one chapter file depends on a \texttt{\textbackslash{}label} or bibliography entry that may live in a different file entirely, and the two are connected only through the auxiliary files LaTeX writes during compilation, not through anything visible in the source itself. This is why a full project typically needs two compilation passes at minimum: the first pass records every label's position and every citation's use, and the second resolves references against that record. A \texttt{\textbackslash{}ref} that renders as a question mark, or a citation that renders as \texttt{{[}?{]}}, almost always means either that a second pass has not yet run or that the file defining that label or entry has not been compiled at all, which happens easily in a multi-file project if \texttt{\textbackslash{}includeonly} is left restricting compilation to a subset of chapters while a reference to another chapter is added.

The underlying requirement, stripped of LaTeX-specific mechanics, is that the pieces of a large document have to agree with each other wherever they touch: a term used in Chapter 12 has to match its definition in Chapter 3, a label referenced in the introduction has to exist somewhere in a later chapter, and the bibliography has to contain every work any chapter cites. None of this is enforced automatically; it is enforced by the discipline of the directory and file conventions in this chapter, and by treating a failed cross-reference or an inconsistent term not as a typo to patch locally but as a sign that two pieces of the document have drifted out of agreement and need to be reconciled.

The remainder of Part II builds on this foundation: version control conventions for a project organized this way, and the build tooling that makes recompiling a project of this size fast enough to work with on a daily basis.
